\origpage{99}{GRAPH MINORS. XIII}
Let $H'$ be a subwall of a wall $H$ in a graph $G$. We say $H'$ is \textit{dividing} if there is no $H$-path in $G$ with one end inside $H'$ and the other end outside, and \textit{non-dividing} otherwise. The following result is closely related to Theorem (6.3) of [21] and its proof is along similar lines (although much easier).

\stmt
(9.2) \textit{For any $l \ge 3$, and $m, p \ge 0$, there exist $k, r \ge 0$ such that the following holds. Let $H$ be a wall with distance function $d$ in a graph $G$, and let $H_1, \dots, H_k$ be non-dividing subwalls of $H$, with}
\begin{equation*}
d(H_i, H_j) \ge r \quad (1 \le i < j \le k).
\end{equation*}
\textit{Then either there is an $(l, m)$-star over $H$ or $G$ has a $K_p$-minor. Moreover, there is an algorithm with running time $\le C\varepsilon(G)$, where $C$ depends only on $l, m, p$ and the height of $H$, which given such subwalls $H_1, \dots, H_k$, finds either an $(l, m)$-star over $H$ or a model of $K_p$ in $G$.}

\textit{Proof.} Let $k, r$ be sufficiently large, in terms of $l, m, p$. (We shall specify later how to choose $k$ and $r$.) Then $H_1, \dots, H_k$ are mutually disjoint.

Let $F$ be a minimal subgraph of $G$ such that for $1 \le i \le k$, $F$ contains an $H$-path from inside $H_i$ to outside $H_i$. Thus $E(F \cap H) = \emptyset$. We can find $F$ in time $O(\varepsilon)$, for fixed $l, m, p, k, r$. \fix{In fact the algorithm needs only a subgraph $F'$ of $G$ with the same property, of bounded size except for divalent vertices, and we find such a subgraph as follows. Let $1 \le i \le k$. Since $H_i$ is non-dividing, there is an $H$-path of $G$ from inside $H_i$ to outside $H_i$, and we find one, $Q_i$ say, in time $O(\varepsilon)$ by breadth-first search from the set of vertices inside $H_i$, using only edges not in $E(H)$ and continuing only from vertices not in $V(H)$, until a vertex of $H$ outside $H_i$ is reached. We build a subgraph $F'$ of $G$ in $k$ steps, starting from the null graph. At step $i$, if no internal vertex of $Q_i$ belongs to $F'$, we add $Q_i$ to $F'$. Otherwise let $z$ be the first internal vertex of $Q_i$ (from its end inside $H_i$) that belongs to $F'$, and add the subpath of $Q_i$ between its end inside $H_i$ and $z$. Then $z \notin V(H)$, since $z$ is an internal vertex of the $H$-path $Q_i$. Let $j < i$ be the first step at which a path containing $z$ was added, and let $A$ be the subpath of this path between its end inside $H_j$ and $z$. The subpath of $Q_i$ between its end inside $H_i$ and $z$ meets $A$ only in $z$, because no vertex of $Q_i$ before $z$ is in $F'$ except possibly the first, which is inside $H_i$, while the only vertex of $A$ in $V(H)$ is its end inside $H_j$. So the union of these two subpaths is an $H$-path of $F'$ from inside $H_i$ to inside $H_j$, and so to outside $H_i$, since $H_i$ and $H_j$ are disjoint. Thus at the end $E(F' \cap H) = \emptyset$, and for $1 \le i \le k$ the graph $F'$ contains an $H$-path from inside $H_i$ to outside $H_i$.

The algorithm uses $F'$ in place of $F$, as follows. By the choice of $z$, each of the $k$ paths added to $F'$ meets the earlier ones only in its ends. Hence the vertices of $H \cup F'$ that are neither branch vertices of $H$ (the vertices that correspond to vertices of the elementary wall) nor ends of these $k$ paths have valency $2$ in $H \cup F'$, and so the size of $H \cup F'$ is bounded, except for divalent vertices, by a function of $k$ and the height of $H$. Moreover $H_1, \dots, H_k$ are non-dividing subwalls of $H$ in the graph $H \cup F'$, and so by the first assertion of (9.2), applied to $H \cup F'$ in place of $G$, there is an $(l, m)$-star over $H$ in $H \cup F'$, or $H \cup F'$ has a $K_p$-minor. Thus at the end of the proof (page 101), where the algorithm searches $H \cup F$, it searches $H \cup F'$ instead. To do so, we call the branch vertices of $H$ and the ends of the $k$ paths \textit{nodes}, and we replace the internal vertices of each path of $H \cup F'$ between two nodes, with no node as an internal vertex, by a single new vertex. This takes time $O(\varepsilon)$. The new graph has bounded size, and in it the images of $H, H_1, \dots, H_k$ are a wall and non-dividing subwalls of it, with the same distances as before. We find an $(l, m)$-star over the new wall or a model of $K_p$ in the new graph by examining all its subgraphs, in time bounded in terms of $l, m, p, k$ and the height of $H$. Replacing each new vertex by the path that it replaced then gives an $(l, m)$-star over $H$ or a model of $K_p$ in $G$, in time $O(\varepsilon)$.}\fixnote{fix:9.2-F-proof-added}{The paper does not say how $F$ is found, nor how $H \cup F$ is searched. The method above is the one used in the formalization.} Since any circuit of $F$ has $\ge 2$ vertices in $V(H)$ (by the minimality of $F$) it follows that we may express $F$ as a union of trees $S_j$ ($1 \le j \le t$), where $S_j \cap S_{j'} \subseteq H$ for all $j \ne j'$, and every vertex of each $S_j$ in $V(H)$ is a leaf of $S_j$, and each leaf of $S_j$ is in $V(H)$.

A vertex $v \in V(F \cap H)$ is \textit{essential} if $v$ is inside $H_i$ for some $i$, and $v$ has valency $1$ in $F$, and no other vertex of $F$ is inside $H_i$.

\hspace*{12pt}(1) \textit{For $1 \le j \le t$, if $e \in E(S_j)$ there exists $i$ with $1 \le i \le k$ such that some leaf $v$ of $S_j$ inside $H_i$ is essential, and every path of $S_j$ between $v$ and another leaf uses $e$.}

\textit{Subproof.} From the minimality of $F$, there exists $i$ ($1 \le i \le k$) such that every $H$-path in $F$ from inside $H_i$ to outside $H_i$ uses $e$. In particular, every such $H$-path is a path of $S_j$, and so for $j' \ne j$, either no leaf of $S_{j'}$ is inside $H_i$ or no leaf of $S_{j'}$ is outside $H_i$. The latter contradicts the minimality of $F$, for then any edge of $S_{j'}$ could be deleted from $F$. Hence for $j' \ne j$, no leaf of $S_{j'}$ is inside $H_i$. Let $v$ be a leaf of $S_j$ inside $H_i$, and suppose that some leaf $u \ne v$ of $S_j$ is either inside $H_i$ or in the perimeter of $H_i$. Let $f$ be the edge of $S_j$ incident with $u$. There is a path of $S_j$ from $v$ to outside $H_i$, and it does not use $f$. Also, for $i' \ne i$, there is a path of $F$ from inside $H_{i'}$ to outside $H_{i'}$, and it can be chosen not to use $f$, by rerouting it to $v$ instead of to $u$ if necessary. But this contradicts the minimality of $F$. Consequently, $v$
